\chapter{Conclusion and Future Work}
\label{ch:conclusion}

\section{Conclusion}

This thesis addressed the industrial need for battery prognostic tools that are
accurate, interpretable, and computationally feasible for deployment on edge
devices within Battery Management Systems (BMS). The literature review established
that existing approaches force a compromise between the accuracy of ``black-box''
deep learning and the interpretability of ``white-box'' physical models, and that
most physics-informed work in the battery domain relies on computationally
expensive PDE-PINNs.

To bridge this gap, the work proposed a \textbf{lightweight hybrid framework} that
regularises a data-driven Gated Recurrent Unit (GRU) with tractable empirical
degradation physics. Two physics mechanisms were implemented and \emph{compared}: a
fixed closed-form prior used as a soft penalty (\textbf{route~A}, using Verhulst
and double-exponential laws) and a \emph{learnable} ODE-residual whose parameters
are trained jointly with the network (\textbf{route~B}). The framework was
evaluated honestly on the full NASA leave-one-cell-out protocol (four folds,
three seeds), a zero-shot cross-dataset transfer to CALCE, a complete physics-design
ablation with a de-confound control, data-scarcity and overfitting stress tests
against two Transformer scales, and a measured Cortex-M4 deployment footprint.

The measured results support a sharper and more defensible conclusion than the
interim ``the hybrid always wins'' narrative:

\begin{itemize}
    \item \textbf{Physics helps --- but the mechanism matters.} The
    \emph{learnable} ODE-residual (route~B) is the robust winner: it attains the
    lowest mean SOH RMSE of all benchmarked models ($0.0307$), is uniform across
    every unseen cell (per-cell RMSE $0.025$--$0.037$), and achieves the best
    physical consistency of any learned model ($0.939$), significantly higher than
    the SOTA-scale Transformer on all 12 folds ($p\approx5\times10^{-4}$). Its
    mean SOH RMSE is $28.3\%$ below that Transformer in-distribution, although
    with only four cells the accuracy gap is not statistically significant
    ($p=0.092$). By contrast, the \emph{fixed} closed-form prior (route~A) helps
    typical cells but actively \emph{hurts} atypical ones --- a prior fit on other
    cells cannot describe a cell such as B0006, which begins above rated capacity
    with heavy regeneration.
    \item \textbf{The Transformer motivation is confirmed and strengthened.} The
    thesis set out from the position that attention models are the accuracy state
    of the art but too heavy for standard BMS hardware. The measurements support
    a stronger claim: under realistic per-fleet data scarcity, scaling the
    Transformer to SOTA size (225k parameters, 13$\times$ the edge-sized variant)
    made it \emph{less} accurate --- classic overfitting, confirmed by scarcity
    and generalisation-gap stress tests --- while the physics-regularised hybrid
    stayed flat from 10\% to 100\% of the training data. On small fleets, the
    physics prior is worth more than $200$k extra attention parameters.
    \item \textbf{Honest baselines.} The LSTM is the strongest purely data-driven
    baseline on unseen NASA cells (RMSE $0.0331$, statistically tied with
    route~B) but shares route~A's weakness on atypical cells, and the light
    Transformer leads the (single-cell, low-confidence) cross-dataset transfer.
    The route-B hybrid is the best \emph{all-round} compromise of accuracy,
    physical plausibility, prognostic calibration (best learned CRA, $0.30$),
    cross-dataset robustness, and footprint --- precisely the balance a
    safety-critical BMS needs.
    \item \textbf{Interpretability with real parameters.} The physics priors were
    fit to real per-cell data (the double-exponential beats the logistic by AIC on
    every cell), and route~B learns its own Verhulst parameters
    ($r\approx-0.80$, $K\approx0.93$) during training, giving the model a
    transparent, physically meaningful degradation law rather than fabricated
    coefficients.
    \item \textbf{Lightweight by measurement, not assertion.} A matched-size
    route-B hybrid (hidden $=32$, $\approx$11k parameters) comes to
    \textbf{11.7~KB} of int8 flash ($4\times$ post-training-quantisation
    estimate of the measured fp32 size) with \textbf{1.25~KB} of measured
    activation RAM ---
    $1/20$ the flash and $1/11$ the latency of the SOTA-scale Transformer ---
    while being \emph{more} accurate ($0.0305$ vs.\ $0.0428$) and retaining the
    full-size model's leave-one-cell-out accuracy. It uses 1.1\% of the
    Cortex-M4 flash budget ($<$256~KB RAM, $<$1~MB flash).
    \item \textbf{Reported transparently.} The homoscedastic uncertainty weighting
    did not have the impact we anticipated on this small, clean problem --- a
    well-chosen fixed weight performed as well or slightly better --- so the fixed
    weight is adopted for the headline models. Similarly, the de-confound ablation
    attributes route~B's raw accuracy largely to its cycle-coordinate input, with
    the ODE residual contributing the (statistically significant) physical
    plausibility. Both findings are stated openly rather than hidden; the
    automatic weighting is expected to earn its place on larger, noisier datasets.
\end{itemize}

In sum, the thesis demonstrates that a \emph{lightweight, learnable
physics-regularised recurrent model} is a viable, trustworthy, and deployable
pathway for the next generation of EV battery prognostics --- accurate where it
matters, physically plausible by construction, robust to unseen cells and data
scarcity, and small enough for on-board inference.

\section{Future Work}

The framework is complete and reproducible end-to-end. Several avenues would
extend its value:

\begin{itemize}
    \item \textbf{Dynamic relaxation of the physics penalty.} The fixed-prior
    route~A degraded atypical cells. A schedule that relaxes the ODE-penalty weight
    in early epochs (letting the network establish empirical convergence before the
    physical constraint tightens) could recover route~A's accuracy on typical cells
    without its failure mode on atypical ones.
    \item \textbf{On-board deployment validation.} The int8 sizes are $4\times$
    quantisation estimates and the Cortex-M4 latency is an analytic MACs/clock
    estimate. Producing a genuinely quantised export (ONNX or TFLite-Micro),
    flashing it to a physical M4 board, and measuring wall-clock inference and
    energy per window would close the deployment loop.
    \item \textbf{Broader chemistries and operating conditions.} Both datasets are
    LCO; validating on LFP/NMC cells and on field-collected partial-cycle data would
    test the framework beyond the cross-load/cross-form-factor transfer studied
    here.
    \item \textbf{More cells for statistical power.} With four NASA cells the
    hybrid's large mean accuracy edge over the SOTA Transformer ($-28.3\%$) falls
    short of significance; evaluating on larger fleets (e.g.\ the Toyota/MIT
    dataset) would settle the accuracy question and re-test the uncertainty
    weighting where it is expected to matter.
    \item \textbf{Per-cell online adaptation.} Because route~B learns its ODE
    parameters, a lightweight online update of those parameters during deployment
    could let a single shipped model specialise to the cell it is actually
    monitoring.
\end{itemize}
